\section{Images and diagrams}

\HomeworkFigure[0.58\linewidth]
  {example-image-a}
  {A sample raster-compatible figure included by the TeX distribution.}
  {A pale rectangle marked with the letter A.}
  {fig:sample-image}

Every figure takes a visual description as its fourth argument, separate from the caption.
The caption is a label for a reader who can see the figure; the description is the figure, in words, for one who cannot.
\Cref{fig:sample-image} carries both.

\HomeworkDiagram[0.42\linewidth]
  {workflow}
  {The source-to-submission workflow, drawn in D2.}
  {A downward flow of four boxes: problem fragments go through lualatex to validation, and validation produces homework.pdf.}
  {fig:workflow}

Diagrams are written in D2 under \filepath{assets/diagrams/} and compiled to PDF by \code{make diagrams}.
The generated PDFs stay committed, so editing a diagram needs D2 installed but compiling the document does not.

A diagram needs its description more than a photograph does, and here for two reasons rather than one.

The first is that the labels are not in the text layer at all.
The pipeline goes through PostScript on its way to PDF, and that hop outlines the glyphs, so \code{pdftotext} returns nothing for a diagram.
That is a deliberate trade: the PostScript route produces a byte-identical PDF on every run, which is what lets the build assert that a committed diagram still matches its source, and the direct route keeps the text but changes bytes between runs.

The second reason would hold even if the words were there.
They would arrive as a loose sequence in drawing order, with no edges, no direction and no shape.
Extractable text is not a description.
